\chapter{Regularized Linear Models}
\label{ch:regularized}

\begin{projectconnection}
Ridge regression is your baseline model throughout the project.
In Phase~1 you build a ridge baseline that appears on \emph{every} chart in Phases~2--5.
If LightGBM doesn't beat ridge on identical features, you haven't learned anything nonlinear---Kozak, Nagel, and Santosh (2020) showed that ridge on principal components matches nonlinear ML methods for SDF estimation across 50 anomaly portfolios.
This chapter teaches why.
\end{projectconnection}

%% ═══════════════════════════════════════════════════════════
\section{Ridge Regression}
\label{sec:ridge}
%% ═══════════════════════════════════════════════════════════

\begin{prereq}[Linear Algebra: Positive-Definite Matrices]
A symmetric matrix $\mathbf{A} \in \R^{p \times p}$ is \emph{positive definite} if $\bx^\top \mathbf{A} \bx > 0$ for all $\bx \neq \mathbf{0}$.
If $\bX^\top \bX$ has eigenvalue $\lambda_{\min} = 0$ (i.e.\ columns are linearly dependent), its inverse does not exist and OLS blows up.
Ridge fixes this by adding $\lambda \mathbf{I}$ to the diagonal, guaranteeing the smallest eigenvalue is at least $\lambda > 0$.
\end{prereq}

\begin{prereq}[Matrix Calculus Refresher]
The gradient of a quadratic $f(\bbeta) = \bbeta^\top \mathbf{A} \bbeta - 2\mathbf{b}^\top \bbeta$ is $\nabla f = 2\mathbf{A}\bbeta - 2\mathbf{b}$.
Setting this to zero gives $\bbeta = \mathbf{A}^{-1}\mathbf{b}$.
We use this twice: once for OLS, once for ridge.
\end{prereq}

\subsection{The Objective}

Ordinary least squares minimizes the residual sum of squares:
\begin{equation}
\hat{\bbeta}_{\mathrm{OLS}} = \arg\min_{\bbeta} \; \| \by - \bX \bbeta \|_2^2
\label{eq:ols-objective}
\end{equation}
where $\by \in \R^n$ is the response vector, $\bX \in \R^{n \times p}$ is the design matrix, and $\bbeta \in \R^p$ is the coefficient vector.
When $p$ is large relative to $n$, or when features are correlated, OLS coefficients become wildly unstable---small perturbations in the data cause large swings in $\hat{\bbeta}$.

Ridge regression (Hoerl and Kennard, 1970) adds an $L_2$ penalty:
\begin{equation}
\hat{\bbeta}_{\mathrm{ridge}} = \arg\min_{\bbeta} \; \underbrace{\| \by - \bX \bbeta \|_2^2}_{\text{fit to data}} + \underbrace{\lambda \| \bbeta \|_2^2}_{\text{shrinkage penalty}}
\label{eq:ridge-objective}
\end{equation}

Term by term:
\begin{itemize}[nosep]
  \item $\| \by - \bX \bbeta \|_2^2 = \sum_{i=1}^n (y_i - \bx_i^\top \bbeta)^2$: sum of squared residuals---measures how well the model fits training data.
  \item $\| \bbeta \|_2^2 = \sum_{j=1}^p \beta_j^2$: the squared $L_2$ norm of the coefficient vector---penalizes large coefficients.
  \item $\lambda \geq 0$: the regularization parameter. At $\lambda = 0$ you get OLS. As $\lambda \to \infty$, all coefficients shrink toward zero.
\end{itemize}

\subsection{Closed-Form Solution}

Expanding the objective and taking the gradient:
\begin{align}
\loss(\bbeta) &= (\by - \bX\bbeta)^\top(\by - \bX\bbeta) + \lambda \bbeta^\top \bbeta \notag \\
\nabla_{\bbeta} \loss &= -2\bX^\top(\by - \bX\bbeta) + 2\lambda \bbeta = \mathbf{0}
\end{align}

Solving for $\bbeta$:
\begin{equation}
\boxed{\hat{\bbeta}_{\mathrm{ridge}} = (\bX^\top \bX + \lambda \mathbf{I})^{-1} \bX^\top \by}
\label{eq:ridge-closed-form}
\end{equation}

Term by term:
\begin{itemize}[nosep]
  \item $\bX^\top \bX \in \R^{p \times p}$: the Gram matrix. Its eigenvalues measure the ``spread'' of data along each direction.
  \item $\lambda \mathbf{I}$: adds $\lambda$ to every eigenvalue of $\bX^\top \bX$. This is what makes the matrix invertible even when $p > n$.
  \item $\bX^\top \by \in \R^p$: the cross-correlation between each feature and the response.
\end{itemize}

\begin{intuition}[What Ridge Actually Does]
Let $\bX = \mathbf{U} \mathbf{D} \mathbf{V}^\top$ be the SVD with singular values $d_1 \geq d_2 \geq \cdots \geq d_p$.
OLS amplifies signal along directions with small $d_j$ (noisy directions get enormous weights).
Ridge replaces $d_j^{-1}$ with $d_j / (d_j^2 + \lambda)$---directions with small singular values get shrunk toward zero.
The noisy directions are precisely the ones that cause OLS to overfit, and ridge damps them out.
\end{intuition}

\subsection{Bias--Variance Tradeoff}

The mean squared error of any estimator decomposes as:
\begin{equation}
\mathrm{MSE}(\hat{\bbeta}) = \underbrace{\| \E[\hat{\bbeta}] - \bbeta \|^2}_{\text{bias}^2} + \underbrace{\tr\!\big(\mathrm{Cov}(\hat{\bbeta})\big)}_{\text{variance}}
\label{eq:bias-variance}
\end{equation}

For ridge, these terms are (using the eigendecomposition $\bX^\top\bX = \mathbf{V}\mathbf{\Lambda}\mathbf{V}^\top$ with eigenvalues $\lambda_1, \ldots, \lambda_p$):
\begin{align}
\text{Bias}^2 &= \sum_{j=1}^p \frac{\lambda^2}{(\lambda_j + \lambda)^2} \alpha_j^2 \label{eq:ridge-bias} \\[4pt]
\text{Variance} &= \sigma^2 \sum_{j=1}^p \frac{\lambda_j}{(\lambda_j + \lambda)^2} \label{eq:ridge-variance}
\end{align}
where $\alpha_j = \mathbf{v}_j^\top \bbeta$ is the projection of the true coefficient vector onto the $j$-th eigenvector and $\sigma^2$ is the noise variance.

The bias term in~\eqref{eq:ridge-bias} increases with $\lambda$: larger penalty means more systematic distortion.
The variance term in~\eqref{eq:ridge-variance} decreases with $\lambda$: regularization stabilizes the estimator.
Hoerl and Kennard (1970) proved the foundational result: there always exists a $\lambda > 0$ for which the ridge estimator has lower MSE than OLS.

\begin{keyresult}[Hoerl--Kennard Theorem (1970)]
For any linear model $\by = \bX\bbeta + \bm{\varepsilon}$ with $\E[\bm{\varepsilon}] = \mathbf{0}$ and $\mathrm{Cov}(\bm{\varepsilon}) = \sigma^2 \mathbf{I}$, there always exists a $\lambda > 0$ such that $\mathrm{MSE}(\hat{\bbeta}_{\mathrm{ridge}}) < \mathrm{MSE}(\hat{\bbeta}_{\mathrm{OLS}})$.
The intuition: OLS variance is $\sigma^2 \sum_j \lambda_j^{-1}$, which can be enormous when any eigenvalue $\lambda_j$ is small.
A small amount of bias buys a large variance reduction.
The optimal $\lambda$ depends on unknowns ($\bbeta$ and $\sigma^2$), so in practice we use cross-validation.
\end{keyresult}

\subsection{Worked Example: Two-Feature Ridge}

\begin{workedexample}{Ridge Shrinkage on Correlated Features}
Consider $n = 100$ observations, $p = 2$ features with correlation $\rho = 0.95$, true $\bbeta = (1, -1)^\top$, noise $\sigma = 2$.

\textbf{OLS result}: Because the features are nearly collinear, the eigenvalues of $\bX^\top\bX$ are approximately $\lambda_1 \approx 195$ and $\lambda_2 \approx 5$.
OLS amplifies noise along the small-eigenvalue direction, giving (say) $\hat{\bbeta}_{\mathrm{OLS}} = (3.2, -3.4)^\top$ with high variance.

\textbf{Ridge at $\lambda = 10$}: The effective eigenvalues become $195 + 10 = 205$ and $5 + 10 = 15$.
The shrinkage ratio along each direction is $\lambda_j / (\lambda_j + \lambda)$: $195/205 = 0.95$ for direction~1 and $5/15 = 0.33$ for direction~2.
Ridge barely touches the well-determined direction but shrinks the noisy direction by $67\%$.
Result: $\hat{\bbeta}_{\mathrm{ridge}} \approx (1.3, -1.2)^\top$---much closer to the truth.

\textbf{Takeaway}: Ridge doesn't shrink all coefficients equally.
It selectively damps the directions where data provides least information.
\end{workedexample}

\subsection{Geometric View: Contour Plot}

Figure~\ref{fig:ridge-contour} shows the ridge solution geometrically. The OLS loss function has elliptical contours (elongated because features are correlated). The $L_2$ constraint $\|\bbeta\|_2^2 \leq t$ defines a circular region. The ridge solution is where the smallest ellipse touches the circle.

\begin{figure}[htbp]
\centering
\begin{tikzpicture}[scale=1.8]
  % Axes
  \draw[->] (-2.3,0) -- (2.8,0) node[right] {$\beta_1$};
  \draw[->] (0,-2.3) -- (0,2.8) node[above] {$\beta_2$};

  % L2 constraint circle
  \draw[thick, defblue, fill=blue!5] (0,0) circle (1.3);
  \node[defblue] at (-0.5, -0.6) {$\|\bbeta\|_2^2 \leq t$};

  % OLS contours (rotated ellipses -- correlated features)
  \begin{scope}[rotate=35]
    \draw[keyorange, thick] (1.4, 0.6) ellipse (0.5 and 1.4);
    \draw[keyorange] (1.4, 0.6) ellipse (0.8 and 1.8);
    \draw[keyorange, dashed] (1.4, 0.6) ellipse (1.1 and 2.2);
  \end{scope}

  % OLS solution
  \fill[warnred] (1.8, 1.8) circle (2pt) node[above right] {$\hat{\bbeta}_{\mathrm{OLS}}$};

  % Ridge solution (where ellipse touches circle)
  \fill[intgreen] (0.75, 1.02) circle (2pt) node[above left] {$\hat{\bbeta}_{\mathrm{ridge}}$};

  % Arrow from OLS to ridge
  \draw[-{Stealth[length=2mm]}, thick, gray, dashed] (1.75, 1.75) -- (0.8, 1.07);
  \node[gray, font=\small] at (1.7, 1.2) {shrinkage};
\end{tikzpicture}
\caption{Ridge regression geometry. The OLS solution (red) lies outside the $L_2$ constraint ball (blue). Ridge finds the point on the ball closest to the OLS optimum (green).}
\label{fig:ridge-contour}
\end{figure}

%% ═══════════════════════════════════════════════════════════
\section{Lasso Regression}
\label{sec:lasso}
%% ═══════════════════════════════════════════════════════════

\begin{prereq}[$L_1$ vs.\ $L_2$ Norms Geometrically]
The $L_2$ ball $\{ \bbeta : \|\bbeta\|_2 \leq t \}$ is a \emph{circle} (sphere in higher dimensions).
The $L_1$ ball $\{ \bbeta : \|\bbeta\|_1 \leq t \}$ is a \emph{diamond} (cross-polytope) with corners on the axes.
The corners are the key: when an elliptical contour meets a diamond, it tends to hit a corner, which means one or more coefficients are exactly zero.
That is why $L_1$ produces sparsity and $L_2$ does not.
\end{prereq}

The Lasso (Tibshirani, 1996) replaces the $L_2$ penalty with an $L_1$ penalty:
\begin{equation}
\hat{\bbeta}_{\mathrm{lasso}} = \arg\min_{\bbeta} \; \| \by - \bX \bbeta \|_2^2 + \lambda \| \bbeta \|_1
\label{eq:lasso-objective}
\end{equation}
where $\| \bbeta \|_1 = \sum_{j=1}^p |\beta_j|$.

Term by term:
\begin{itemize}[nosep]
  \item The fit term $\| \by - \bX\bbeta \|_2^2$ is identical to ridge.
  \item $\| \bbeta \|_1$: the sum of absolute values. Unlike the smooth $L_2$ norm, this has a kink at zero, which is what forces coefficients to hit exactly zero.
  \item $\lambda$: same role as in ridge, but now controls how many coefficients survive (larger $\lambda$ = fewer nonzero coefficients).
\end{itemize}

\subsection{No Closed Form---Coordinate Descent}

Because the absolute value $|\beta_j|$ is not differentiable at zero, there is no closed-form solution.
The standard algorithm is \textbf{coordinate descent}: cycle through coefficients one at a time, holding all others fixed.
The update for coefficient $j$ uses the \emph{soft-thresholding operator}:
\begin{equation}
\hat{\beta}_j \leftarrow S\!\left(\frac{1}{n}\sum_{i=1}^n x_{ij} r_i^{(-j)}, \; \frac{\lambda}{2n}\right), \qquad S(z, \gamma) = \operatorname{sign}(z)\max(|z| - \gamma, 0)
\label{eq:soft-threshold}
\end{equation}
where $r_i^{(-j)} = y_i - \sum_{k \neq j} x_{ik}\hat{\beta}_k$ is the partial residual.
The key insight: when the partial correlation $|z|$ is smaller than the threshold $\gamma = \lambda/(2n)$, the coefficient is set to exactly zero.
This is the mechanism that produces sparsity---weak signals get zeroed out.

An alternative is \textbf{LARS} (Least Angle Regression), which traces the full regularization path efficiently in $O(np \min(n,p))$ operations (Efron et al., 2004).
In practice, coordinate descent (as implemented in \texttt{sklearn}) is faster for large-scale problems.

\subsection{Why $L_1$ Produces Sparsity}

The geometric argument is simple (Figure~\ref{fig:lasso-contour}).
The Lasso constraint region is a diamond with corners on the coordinate axes.
When the OLS contour ellipses touch this diamond, they almost always hit a corner first---and a corner means $\beta_j = 0$ for at least one $j$.
Ridge's circle has no corners, so the tangent point generically has all coefficients nonzero.

\begin{figure}[htbp]
\centering
\begin{tikzpicture}[scale=1.8]
  % Axes
  \draw[->] (-2.3,0) -- (2.8,0) node[right] {$\beta_1$};
  \draw[->] (0,-2.3) -- (0,2.8) node[above] {$\beta_2$};

  % L1 constraint diamond
  \draw[thick, defblue, fill=blue!5]
    (1.3,0) -- (0,1.3) -- (-1.3,0) -- (0,-1.3) -- cycle;
  \node[defblue] at (-0.8, -0.6) {$\|\bbeta\|_1 \leq t$};

  % OLS contours (rotated ellipses)
  \begin{scope}[rotate=25]
    \draw[keyorange, thick] (1.2, 0.8) ellipse (0.4 and 1.2);
    \draw[keyorange] (1.2, 0.8) ellipse (0.7 and 1.6);
    \draw[keyorange, dashed] (1.2, 0.8) ellipse (1.0 and 2.0);
  \end{scope}

  % OLS solution
  \fill[warnred] (1.8, 1.5) circle (2pt) node[above right] {$\hat{\bbeta}_{\mathrm{OLS}}$};

  % Lasso solution (on a corner of the diamond -- sparse!)
  \fill[intgreen] (0, 1.3) circle (2pt) node[above right] {$\hat{\bbeta}_{\mathrm{lasso}}$};

  % Arrow from OLS to lasso
  \draw[-{Stealth[length=2mm]}, thick, gray, dashed] (1.75, 1.45) -- (0.08, 1.32);
  \node[gray, font=\small] at (1.5, 0.9) {shrinkage};
  \node[intgreen, font=\small] at (0.8, 1.5) {$\beta_1 = 0$ (sparse!)};
\end{tikzpicture}
\caption{Lasso geometry. The diamond constraint has corners on the axes. The OLS contour hits the corner at $\beta_1 = 0$, producing a sparse solution.}
\label{fig:lasso-contour}
\end{figure}

\subsection{Ridge vs.\ Lasso: Comparison}

\begin{center}
\begin{tabular}{lcc}
\toprule
\textbf{Property} & \textbf{Ridge ($L_2$)} & \textbf{Lasso ($L_1$)} \\
\midrule
Penalty & $\lambda \sum_j \beta_j^2$ & $\lambda \sum_j |\beta_j|$ \\
Sparsity & No (all $\beta_j \neq 0$) & Yes (many $\beta_j = 0$) \\
Closed form & Yes & No \\
Correlated features & Shrinks together & Picks one arbitrarily \\
Constraint geometry & Circle & Diamond \\
Best when & Many small effects & Few large effects \\
\bottomrule
\end{tabular}
\end{center}

The correlated-features row matters for your project: risk-cube features (VaR utilization, factor concentration, scenario P\&L) will be correlated.
Lasso would arbitrarily drop one of two correlated signals.
Ridge keeps both, shrunk.

%% ═══════════════════════════════════════════════════════════
\section{Elastic Net}
\label{sec:elastic-net}
%% ═══════════════════════════════════════════════════════════

Elastic net (Zou and Hastie, 2005) combines $L_1$ and $L_2$:
\begin{equation}
\hat{\bbeta}_{\mathrm{enet}} = \arg\min_{\bbeta} \; \| \by - \bX\bbeta \|_2^2 + \lambda \Big[ \alpha \|\bbeta\|_1 + \frac{1-\alpha}{2} \|\bbeta\|_2^2 \Big]
\label{eq:elastic-net}
\end{equation}

Term by term:
\begin{itemize}[nosep]
  \item $\alpha \in [0,1]$: the mixing parameter. At $\alpha = 1$ you get lasso; at $\alpha = 0$ you get ridge.
  \item $\lambda$: overall regularization strength, same role as before.
  \item The $L_1$ component produces sparsity.
  \item The $L_2$ component handles correlated features by grouping them.
\end{itemize}

\begin{keyidea}[When to Use Elastic Net]
Use elastic net when you suspect a sparse model but have correlated features.
The $L_2$ term ensures that if features $j$ and $k$ are highly correlated, elastic net includes both (with shrunken coefficients) rather than arbitrarily dropping one.
Zou and Hastie (2005) called this the ``grouping effect.''

In practice, \texttt{sklearn}'s \texttt{ElasticNetCV} selects both $\alpha$ and $\lambda$ by cross-validation.
A common default grid: $\alpha \in \{0.1, 0.5, 0.7, 0.9, 0.95, 1.0\}$.
\end{keyidea}

For your project, the choice is usually ridge (baseline) or LightGBM (nonlinear).
Elastic net sits between them---useful if you want a sparse linear model, but ridge is the cleaner baseline because it doesn't introduce a variable-selection step that could interact with your validation procedure.

\begin{workedexample}{Elastic Net on Correlated Risk Features}
Suppose you have three features: VaR utilization ($x_1$), VaR utilization lagged one day ($x_2$), and factor Herfindahl ($x_3$).
Features $x_1$ and $x_2$ have correlation $\rho = 0.98$; $x_3$ is independent.

\textbf{Lasso} ($\alpha = 1$): picks one of $x_1, x_2$ arbitrarily, zeros the other. You lose the lag structure.

\textbf{Ridge} ($\alpha = 0$): keeps all three with shrunken coefficients.
$\hat{\bbeta} \approx (0.3, 0.25, 0.4)^\top$.
No sparsity, but the lag information is preserved.

\textbf{Elastic net} ($\alpha = 0.5$): groups $x_1$ and $x_2$ together (both get similar nonzero coefficients), while also providing some sparsity pressure.
If $x_3$ were pure noise, elastic net might zero it out while keeping both $x_1$ and $x_2$.
\end{workedexample}

%% ═══════════════════════════════════════════════════════════
\section{Ridge on Principal Components}
\label{sec:ridge-pca}
%% ═══════════════════════════════════════════════════════════

\begin{prereq}[Principal Component Analysis]
PCA decomposes $\bX = \mathbf{U}\mathbf{D}\mathbf{V}^\top$ and projects onto the columns of $\mathbf{V}$ (the right singular vectors).
The first PC captures the direction of maximum variance, the second captures the most variance orthogonal to the first, and so on.
Projecting onto the top $k$ PCs gives $\tilde{\bX} = \bX \mathbf{V}_k \in \R^{n \times k}$, reducing dimensionality from $p$ to $k$.
\end{prereq}

Ridge regression on principal components is a two-step procedure:
\begin{enumerate}[nosep]
  \item Compute principal components $\mathbf{Z} = \bX \mathbf{V}_k$ (keep top $k$ PCs).
  \item Run ridge on $\mathbf{Z}$: $\hat{\bm{\gamma}} = (\mathbf{Z}^\top\mathbf{Z} + \lambda \mathbf{I})^{-1}\mathbf{Z}^\top\by$.
  \item Map back: $\hat{\bbeta} = \mathbf{V}_k \hat{\bm{\gamma}}$.
\end{enumerate}

This is closely related to plain ridge---in fact, ridge regression on all features with penalty $\lambda$ is equivalent to shrinking PCs with shrinkage factor $d_j^2 / (d_j^2 + \lambda)$, where $d_j$ is the $j$-th singular value.
Low-variance PCs (small $d_j$) get shrunk most aggressively.
The difference with explicit PC regression is that you choose $k$ and then apply ridge, giving you two knobs instead of one.

\begin{keyresult}[Kozak, Nagel, Santosh (2020, JFE)]
Using 50 anomaly-sorted portfolios from the cross-section of stock returns, the authors show:
\begin{itemize}[nosep]
  \item A PC-sparse SDF (i.e.\ ridge regression on the top principal components of characteristics-based factor returns) achieves high out-of-sample $R^2$ with as few as \textbf{four principal components}.
  \item No similarly sparse SDF based on the \emph{primitive} characteristics-based factors can match this out-of-sample performance---characteristics-sparse models require far more factors and still underperform.
  \item The data ``calls for substantial $L_2$-shrinkage but essentially no sparsity''---ridge dominates lasso for SDF estimation.
  \item Unregularized models (OLS on all 50 factors) produce out-of-sample $R^2$ substantially below zero.
\end{itemize}
\textbf{Implication}: Ridge on PCs is a strong default that nonlinear methods must beat. Won the 2020 Fama--DFA Prize for best paper in the JFE.
\end{keyresult}

Why does this matter for your project? You are building signals from a high-dimensional set of risk-cube features (VaR utilization, factor Herfindahl, scenario P\&L ranks, etc.). If ridge on the PCs of these features already captures the predictable variation in returns, adding LightGBM buys you nothing---you've just overfit. Ridge on PCs is the acid test.

\begin{definition}[Effective Degrees of Freedom for Ridge]
The effective degrees of freedom of the ridge estimator is:
\begin{equation}
\mathrm{df}(\lambda) = \tr(\mathbf{H}_\lambda) = \sum_{j=1}^p \frac{d_j^2}{d_j^2 + \lambda}
\label{eq:effective-df}
\end{equation}
At $\lambda = 0$: $\mathrm{df} = p$ (full OLS). As $\lambda \to \infty$: $\mathrm{df} \to 0$.
This is a continuous analog of ``number of parameters''---it measures model complexity under regularization and appears in information criteria (AIC, BIC) adapted for ridge.
\end{definition}

%% ═══════════════════════════════════════════════════════════
\section{Regularization Path and $\lambda$ Selection}
\label{sec:lambda-selection}
%% ═══════════════════════════════════════════════════════════

\begin{prereq}[Cross-Validation Basics]
$K$-fold CV splits data into $K$ folds, trains on $K{-}1$ folds, and evaluates on the held-out fold.
Repeat for each fold and average.
The value of $\lambda$ that minimizes average held-out error is selected.
Common choice: $K = 5$ or $K = 10$.
For time series, use \emph{purged} $K$-fold CV (Chapter~\ref{ch:validation}) to avoid look-ahead bias.
\end{prereq}

\subsection{The Regularization Path}

As $\lambda$ increases from $0$ to $\infty$, each coefficient traces a path from its OLS value to zero.
Plotting these paths reveals which features are ``most important'' (persist longest) and which are noise (shrink to zero quickly).

For lasso, this is called the ``regularization path'' and can be computed exactly using LARS.
For ridge, the path is $\hat{\beta}_j(\lambda) = \sum_k v_{jk} \cdot d_k \cdot u_k^\top \by / (d_k^2 + \lambda)$, where the $v_{jk}$, $d_k$, $u_k$ come from the SVD.

\subsection{$\lambda$ Selection: Cross-Validation}

\textbf{For ridge}: use generalized cross-validation (GCV), a leave-one-out shortcut:
\begin{equation}
\mathrm{GCV}(\lambda) = \frac{1}{n} \sum_{i=1}^n \left( \frac{y_i - \hat{y}_i(\lambda)}{1 - h_{ii}(\lambda)} \right)^2
\label{eq:gcv}
\end{equation}
where $h_{ii}(\lambda)$ is the $i$-th diagonal of the hat matrix $\mathbf{H}_\lambda = \bX(\bX^\top\bX + \lambda\mathbf{I})^{-1}\bX^\top$.
GCV computes leave-one-out CV \emph{exactly} in $O(np^2)$---no need to refit $n$ times.

\textbf{For lasso/elastic net}: no LOO shortcut exists. Use $K$-fold CV. The \texttt{sklearn} classes \texttt{LassoCV} and \texttt{ElasticNetCV} handle this automatically.

\textbf{Common practice}: search $\lambda$ on a log-scale grid, e.g.\ $\lambda \in \{10^{-4}, 10^{-3}, \ldots, 10^{4}\}$, then refine around the minimum.
The ``one-standard-error rule'' is also common: instead of picking the $\lambda$ that minimizes CV error, pick the largest $\lambda$ whose CV error is within one standard error of the minimum.
This yields a simpler (more regularized) model that performs nearly as well.

\begin{warning}[Time Series: Don't Use Standard $K$-Fold]
Standard $K$-fold CV assumes observations are i.i.d.
Financial returns are serially correlated: if observation $t$ is in the training set and $t{+}1$ is in the validation set, the model has seen ``almost the answer.''
Use \textbf{purged $K$-fold CV with embargo} (Chapter~\ref{ch:validation}): remove observations within a buffer around each test fold.
Failing to do this inflates your $R^2$ and gives you false confidence.
\end{warning}

%% ═══════════════════════════════════════════════════════════
\section{Why Ridge Is the Baseline in Finance}
\label{sec:ridge-baseline}
%% ═══════════════════════════════════════════════════════════

Return prediction lives in a brutally low signal-to-noise regime.

\begin{keyidea}[The Low-SNR Reality of Return Prediction]
Monthly stock return prediction typically yields out-of-sample $R^2$ values of $0.1\%$--$0.5\%$.
Gu, Kelly, and Xiu (2020, RFS) report:
\begin{itemize}[nosep]
  \item Elastic net: monthly OOS $R^2 \approx 0.11\%$ (using 94 stock characteristics plus interactions, totaling 900+ predictors, 1957--2016 sample)
  \item PCR / PLS: $\approx 0.26$--$0.27\%$
  \item Best neural network (NN3): $\approx 0.40\%$
\end{itemize}
These are \emph{tiny} numbers.
An $R^2$ of $0.40\%$ means the model explains less than half of one percent of monthly return variation---yet it is enough to generate an annualized Sharpe ratio of 1.35 for a value-weighted long-short portfolio.
In finance, tiny predictability = real money.
\end{keyidea}

In this regime, the gap between ridge and more complex methods is often small.
The neural network's $0.40\%$ is roughly $4\times$ the elastic net's $0.11\%$, but this gap narrows on cleaner (lower-dimensional) feature sets.
When you're working with 10--20 risk-cube features rather than 900+ stock-level characteristics, ridge may well be sufficient.

\begin{projectconnection}[Phase 2: Your OOS $R^2$ Will Be Small]
With $\sim$1,250 daily observations (5 years) and 10--20 risk-cube features, expect OOS $R^2$ in the range of $0.5$--$3\%$ for daily return prediction---if the signal exists at all.
Do not be discouraged by these numbers.
The question is not ``is $R^2$ large?'' but ``does ridge $R^2 > 0$ after purged CV, and does LightGBM beat it?''
A daily $R^2$ of $1\%$ with stable IC across folds and a Deflated Sharpe Ratio above zero is a publishable result in this domain.
\end{projectconnection}

\begin{keyidea}[Ridge as an Integrity Check]
In your project, ridge serves two purposes:
\begin{enumerate}[nosep]
  \item \textbf{Baseline}: If LightGBM's IC, Sharpe, or $R^2$ doesn't exceed ridge on identical features, you've just overfit noise. The nonlinear model learned nothing that a simple linear combination couldn't capture.
  \item \textbf{Presentation}: Showing ridge alongside LightGBM on every chart signals intellectual honesty. It demonstrates you're testing whether nonlinearity exists, not assuming it.
\end{enumerate}
This is not optional---it's what separates credible quant research from curve-fitting.
\end{keyidea}

\begin{warning}[The Missing Baseline Problem]
Many published ML-for-finance papers report complex model results without a ridge (or OLS) baseline on the same features.
Without this comparison, you cannot distinguish three very different situations:
\begin{enumerate}[nosep]
  \item \textbf{Genuine nonlinear signal}: GBM $\gg$ ridge $>$ zero---the data has real nonlinear structure.
  \item \textbf{Linear signal dressed up}: GBM $\approx$ ridge $>$ zero---the signal is real but linear. GBM wasted compute.
  \item \textbf{Overfitting}: GBM $>$ zero, ridge $\leq$ zero---GBM memorized noise.
\end{enumerate}
Only scenario 1 justifies a nonlinear model.
The Kozak--Nagel--Santosh result warns that scenario 2 is common in asset pricing: the SDF is approximately linear in a small number of PCs.
\end{warning}

\subsection{The McLean--Pontiff Decay Warning}

Even signals that are genuinely predictive in-sample decay after publication.
McLean and Pontiff (2016, JF) examined 97 published stock-return predictors and found:
\begin{itemize}[nosep]
  \item Returns were \textbf{26\% lower} out-of-sample (before publication), suggesting data-mining effects.
  \item Returns were \textbf{58\% lower} post-publication, with an estimated \textbf{32\%} attributable to publication-informed trading.
\end{itemize}
This means even a real signal loses more than half its potency once traders know about it.
Your internal risk-cube signals have an advantage here: they are proprietary and not subject to publication-informed arbitrage.
But the 26\% out-of-sample decay is a data-mining effect that applies everywhere---which is why Deflated Sharpe Ratio (Chapter~\ref{ch:validation}) is non-negotiable.

%% ═══════════════════════════════════════════════════════════
\section{Implementation Notes}
\label{sec:ridge-implementation}
%% ═══════════════════════════════════════════════════════════

Practical points for your codebase:

\begin{enumerate}[nosep]
  \item \textbf{Standardize features first}. Ridge penalizes $\|\bbeta\|_2^2$, which is scale-dependent. If feature~1 is in basis points and feature~2 is a ratio, the penalty is asymmetric. Always standardize to zero mean and unit variance before fitting.

  \item \textbf{\texttt{sklearn} API}:
    \begin{itemize}[nosep]
      \item \texttt{Ridge(alpha=lambda)}: single $\lambda$.
      \item \texttt{RidgeCV(alphas=[...])}: built-in GCV for $\lambda$ selection.
      \item \texttt{Lasso}, \texttt{LassoCV}, \texttt{ElasticNet}, \texttt{ElasticNetCV}: analogous.
    \end{itemize}
    Note: \texttt{sklearn} uses \texttt{alpha} where this chapter uses $\lambda$.

  \item \textbf{Intercept handling}. By default, \texttt{sklearn} fits an unpenalized intercept (mean of $\by$). This is correct---you don't want to shrink the intercept toward zero.

  \item \textbf{$\lambda$ grid}. Start with \texttt{np.logspace(-4, 4, 50)}. If the optimal $\lambda$ hits the boundary, extend the grid.

  \item \textbf{For your project} (\texttt{src/validation/baseline.py}): wrap \texttt{RidgeCV} with purged $K$-fold CV from Chapter~\ref{ch:validation}, not \texttt{sklearn}'s default GCV (which assumes i.i.d.\ data).
\end{enumerate}

%% ═══════════════════════════════════════════════════════════
\section{Summary}
\label{sec:regularized-summary}
%% ═══════════════════════════════════════════════════════════

\begin{center}
\begin{tabular}{p{4cm} p{9.5cm}}
\toprule
\textbf{Concept} & \textbf{Key Takeaway} \\
\midrule
Ridge ($L_2$) & Closed-form: $(\bX^\top\bX + \lambda\mathbf{I})^{-1}\bX^\top\by$. Shrinks all coefficients, never zeros them out. Always exists a $\lambda$ that beats OLS in MSE (Hoerl--Kennard, 1970). \\[6pt]
Lasso ($L_1$) & Produces sparse models ($\beta_j = 0$). No closed form. Arbitrarily picks one of correlated features. \\[6pt]
Elastic net & Combines $L_1 + L_2$. Grouping effect handles correlated features. Two hyperparameters ($\lambda, \alpha$). \\[6pt]
Ridge on PCs & 4 PCs suffice for SDF estimation; ridge on PCs matches nonlinear ML (Kozak--Nagel--Santosh, 2020). \\[6pt]
$\lambda$ selection & GCV for ridge (LOO shortcut), $K$-fold CV for lasso/elastic net. \textbf{Use purged CV for time series}. \\[6pt]
Baseline role & Ridge is the bar GBM must clear. Three scenarios: genuine nonlinear signal, linear signal dressed up, or overfitting. Only scenario 1 justifies nonlinear models. \\[6pt]
Low-SNR reality & Monthly OOS $R^2 \sim 0.1$--$0.4\%$ for stock returns (Gu--Kelly--Xiu, 2020). Published signals decay 58\% post-publication (McLean--Pontiff, 2016). \\
\bottomrule
\end{tabular}
\end{center}

\bigskip
\noindent\textbf{What's next}: Chapter~\ref{ch:trees} covers tree-based methods (random forests, gradient boosting, LightGBM)---the models that must beat the ridge baseline you've just learned to build.
